% Generated draft: requires mathematical and source review.
\documentclass[11pt]{article}
\usepackage[margin=25mm]{geometry}
\usepackage{fontspec}
\setmainfont{Latin Modern Roman}
\usepackage{amsmath,amssymb,mathtools,unicode-math}
\setmathfont{Latin Modern Math}
\usepackage{enumitem,needspace,xurl,hyperref}
\setlength{\parindent}{0pt}
\setlength{\parskip}{5pt}
\setlength{\emergencystretch}{3em}
\newcommand{\N}{\mathbb{N}}
\newcommand{\Z}{\mathbb{Z}}
\newcommand{\Q}{\mathbb{Q}}
\newcommand{\R}{\mathbb{R}}
\newcommand{\C}{\mathbb{C}}
\newcommand{\F}{\mathbb{F}}
\newcommand{\A}{\mathbb{A}}
\newcommand{\PP}{\mathbb{P}}
\newcommand{\E}{\mathbb{E}}
\newcommand{\eps}{\varepsilon}
\newcommand{\dd}{\,\mathrm{d}}
\begin{document}
\section*{UnsolvedMath ID: 30004166}
\textbf{Source problem code:} catalogue entry
\subsection{Definitions and mathematical statement}
Let \(K\) be the function field in the supplied problem, let \(B\) denote its base, and let \(X\) be a smooth projective variety over \(K\). The supplied record does not specify the hypotheses on \(B\), the field of constants, or the dimension of \(X\). Define the rational Picard group by
\[
\operatorname{Pic}(B)_{\Q}:=\operatorname{Pic}(B)\otimes_{\Z}\Q,
\]
where \(\operatorname{Pic}(B)\) is the group of isomorphism classes of line bundles on \(B\), with tensor product as its operation. A cycle is homologically trivial relative to a chosen cycle class map if its image under that map is zero; the relevant cycle class map and cycle groups are not specified in the supplied record.

\textbf{Question.} For homologically trivial cycles on \(X\), does the constructed generalized Beilinson \(\acute{\mathrm{e}}\)tale height pairing arise from a motivic pairing with values in \(\operatorname{Pic}(B)_{\Q}\), with compatibility under \(\acute{\mathrm{e}}\)tale realization recovering the given height pairing?

Here ``arise from a motivic pairing'' retains the source's intended motivic requirement. The record does not specify the motivic category, the two inputs of the pairing, their cycle codimensions or equivalence relations, the coefficient and realization conventions, the target of the constructed \(\acute{\mathrm{e}}\)tale pairing, or the comparison maps. Thus this is a statement of the supplied existence question; a fully specified quantified formulation requires those missing construction data. No particular choices of these data, or additional hypotheses on \(B\), are asserted.
\subsection{Short English statement}
The question asks whether the generalized Beilinson \(\acute{\mathrm{e}}\)tale height pairing for homologically trivial cycles on a smooth projective variety over a function field comes from a motivic pairing valued in the rational Picard group of the base. The supplied record leaves the precise cycle groups, base hypotheses, and realization compatibility unspecified.
\subsection{Sources}
\begin{itemize}
\item Local UnsolvedMath registry, permanent ID \(30004166\), title ``Motivic Origins of \(\acute{\mathrm{E}}\)tale Height Pairings,'' Oberwolfach Reports Open Problems collection, exact source problem code \texttt{catalogue entry}. The supplied proposed year is 2019; no proposer is named.
\item Supplied source citation: \emph{Algebraic K-theory (2020)}, Oberwolfach Reports, DOI 10.4171/owr/2019/29, pp.\ 1737--1790. Source URL and literature-source label ``OWR primary report'': \url{https://doi.org/10.4171/owr/2019/29}. The supplied background describes these reports as documenting mathematical workshops and their open-problem discussions.
\item Supplied literature-source label ``generalized Beilinson geometric height pairing'': \url{https://doi.org/10.4171/DM/X15}. No authors or further bibliographic details are supplied for this reference.
\item Literature assessment checked at \texttt{2026-08-22}: the supplied curation labels the problem open and states that the generalized Beilinson height pairing is constructed and has proved realizations, while its asserted motivic origin remains explicitly conjectural and is described by the authors as speculative. This is supplied dated curation, not an independently verified assessment. The background explicitly requests independent verification of time-sensitive status and cited claims.
\item The registry labels the statement \texttt{corrected\_verified} and reports a manual semantic, typographic, and LaTeX review. These are supplied registry assertions, not independent verification of the formulation or its mathematical completeness.
\end{itemize}
\end{document}
